% Ported from sections/01_methods.md (v6, APPROVED 2026-09-18).
% Final-manuscript Sec.~II. Floats: Table 1, Table 2, Figure 1 (TikZ).
\section{Methods}
\label{sec:methods}

\subsection{Interface models}
\label{sec:models}

We study the metal-on-oxide interface relevant to magnetic tunnel junctions: a thin metal film
deposited on the MgO(001) surface. Two compositions were considered, differing only in whether boron
is present in the film (Table~\ref{tab:models})---the pair that isolates the effect of the added
metalloid in a single host.

\begin{table}[t]
\small
\caption{\label{tab:models}The two interface models: film constitution and total atom count.}
\begin{tabular*}{\textwidth}{@{\extracolsep{\fill}}lllc}
\toprule
Model & Composition & Film species & $N$ atoms \\
\midrule
Fe/MgO & Fe$_{25}$Mg$_{25}$O$_{25}$ & Fe & 75 \\
Fe-B/MgO & B$_3$Fe$_{25}$Mg$_{25}$O$_{25}$ & Fe $+$ B & 78 \\
\bottomrule
\end{tabular*}
\end{table}

The substrate is a single MgO(001) layer and the film a single metal layer, stacked along $z$ with
20~\AA{} of vacuum and periodic boundaries in-plane. The in-plane lattice of the simulation
cell is the DFT-optimised Fe lattice constant ($a_\mathrm{Fe} = 2.870$~\AA), used in a $(5\times5\times1)$
supercell throughout. The MgO layer is built epitaxially onto that geometry: the substrate is
assembled on the Fe(001) surface cell with the oxygen sublattice placed directly above the metal
sites (a $\sqrt{2}$ rotation relative to the bulk MgO cell), so the MgO in-plane parameter is forced
to 2.870~\AA. Compared with the bulk value ($a_\mathrm{MgO}/\sqrt{2} = 2.978$~\AA) this means the
MgO is compressed in-plane by 3.6\,\%, i.e.\ the mismatch is accommodated by the substrate,
not by the film---the Fe film remains at its own equilibrium lattice constant and is not strained
in-plane. The substrate is then held fixed in that compressed state and only the film is
allowed to relax.

\subsection{Biased global optimisation}
\label{sec:bias}

Structures were sampled with the Atomistic Global Optimization X (AGOX) package \cite{agox2020},
using a global optimisation with first-principles energy expressions (GOFEE) search
\cite{gofee2017}: a Gaussian-process regression (GPR) surrogate with an Oganov-type fingerprint
descriptor \cite{oganov2011} and a repulsive prior, coupled to a lower confidence bound (LCB)
acquisition function ($\kappa = 2$). Each search ran 100 iterations.

\subsubsection{Three-phase biased exploration} Candidate generation is not uniform over the run but
proceeds in three phases selected by the iteration counter $i$, mixing generators of different
perturbation scale. Both models combine a small-scale heterostructure-aware randomiser
(rattle amplitude 1.5, all film atoms) with a large-scale rattle generator (half the film
atoms, rattle amplitude 2.3); the schedule is identical for the two models and is given in
Table~\ref{tab:schedule}.

\begin{table}[t]
\small
\caption{\label{tab:schedule}Candidate-generation schedule, common to both models: a small-scale
heterostructure-aware randomiser and a large-scale rattle generator. The number of candidates
generated per phase is $N$.}
\begin{tabular*}{\textwidth}{@{\extracolsep{\fill}}lccc c}
\toprule
Phase & Iterations & Small-scale & Large-scale & Total $N$ \\
\midrule
\textbf{I}   & $0 \le i < 10$   & 20 & 0  & 20 \\
\textbf{II}  & $10 \le i < 25$  & 10 & 10 & 20 \\
\textbf{III} & $25 \le i$       & 0  & 20 & 20 \\
\bottomrule
\end{tabular*}
\end{table}

The per-phase candidate total is 20; the two models differ only in whether the film contains boron,
not in how candidates are generated. The resulting workflow is summarised in Fig.~\ref{fig:loop}.

\begin{figure}[t]
\centering
\begin{tikzpicture}[
  font=\footnotesize,
  >={Stealth[length=1.8mm]},
  node distance=5mm and 7mm,
  box/.style={draw, rounded corners=1pt, align=center, inner sep=3pt},
  dec/.style={draw, diamond, aspect=2.4, align=center, inner sep=0pt},
]
\node[box] (init) {Initialize the search\\ ($i = 0$)};
\node[dec, below=of init] (d1) {$0 \le i < 10$};
\node[dec, below=of d1] (d2) {$10 \le i < 25$};
\node[dec, below=of d2] (d3) {$25 \le i$};

\node[box, right=13mm of d1] (p1) {Phase I\\ Generate $N$ candidates\\ (small-scale)};
\node[box, right=13mm of d2] (p2) {Phase II\\ Generate $N$ candidates\\ (small $+$ large)};
\node[box, right=13mm of d3] (p3) {Phase III\\ Generate $N$ candidates\\ (large-scale)};

\begin{scope}[on background layer]
\node[draw, dashed, rounded corners=2pt, fit=(p1)(p2)(p3), inner sep=5pt] (BE) {};
\end{scope}
\node[above=0pt of BE.north, font=\footnotesize\bfseries] {Biased Exploration};

\node[box, below=10mm of d3] (gpr) {GPR surrogate\\ optimization};
\node[box, below=of gpr] (lcb) {LCB: choose $M$ candidates};
\node[box, below=of lcb] (dft) {DFT evaluate\\ $M$ candidates};
\node[box, below=of dft] (db) {Database};

\draw[->] (init) -- (d1);
\draw[->] (d1) -- node[left,font=\scriptsize]{N} (d2);
\draw[->] (d2) -- node[left,font=\scriptsize]{N} (d3);
\draw[->] (d1) -- node[above,font=\scriptsize]{Y} (p1);
\draw[->] (d2) -- node[above,font=\scriptsize]{Y} (p2);
\draw[->] (d3) -- node[above,font=\scriptsize]{Y} (p3);
\draw[->] (p1) |- (gpr);
\draw[->] (p2) |- (gpr);
\draw[->] (p3) -- (gpr);
\draw[->] (gpr) -- (lcb);
\draw[->] (lcb) -- (dft);
\draw[->] (dft) -- node[right,font=\scriptsize]{store} (db);
\draw[->] (db.west) -- ++(-7mm,0) |- node[pos=0.25,left,font=\scriptsize]{$i+1$} (d1.west);
\end{tikzpicture}
\caption{\label{fig:loop}The biased-exploration loop. Each iteration selects a phase from the
counter $i$, generates $N$ candidates with that phase's generator mix (Table~\ref{tab:schedule}),
optimises them against the GPR surrogate, ranks them with the LCB acquisition, evaluates the
selected candidates with DFT and stores them in the database; the counter then advances and the
phase is re-selected.}
\end{figure}

The early phases favour the small-scale (structure-aware) generator, which preserves the
layer-resolved character of the flat reference basin, while the later phases shift to the
large-scale rattle generator for broader exploration.

\subsubsection{Bias} The exploration is deliberately biased: the search is seeded from a
flat metal layer---the reference film geometry, a single-monolayer-thick film rather than a
random three-dimensional distribution of metal atoms---so the sampled database is a mixture of the
flat basin and any lower-lying basin the search finds. The composition of the reference layer
follows the target stoichiometry (Sec.~\ref{sec:models}): a pure Fe layer for Fe/MgO, and the same Fe
layer with B added above the surface for Fe-B/MgO. The films are single-layer in both models, and the
two reference layers differ only by the boron. Each model was run from multiple independent random
seeds, and each seed constitutes one independent search.

\subsubsection{Only completed searches are used} Every number in this paper comes from searches that ran
the full 100-iteration budget: 13 completed searches for Fe/MgO and 6 for Fe-B/MgO, 19 in
total. Two further searches were started and are \emph{not} used, and we record them here rather than
drop them silently---one Fe/MgO search reached only iteration 37, and one Fe-B/MgO search produced no
stored structures at all. A stopped search has had less search time than a completed one, so its best
energy is systematically worse; and because such searches are not distributed evenly across the two
models, including them would bias the comparison between them---as the method-parameter study in the
Supplementary Material also shows, where searches that stopped early inflate the mean of one setting
and depress another (Sec.~S3 of the Supplementary Material). The same criterion is applied without
exception throughout the analysis, and it is decided from the iteration count recorded for
each search, not from how a search is labelled: a search that looks complete can have stopped early,
and one whose stored structures are missing is treated the same way. In the two models used here the
excluded searches happen to be identifiable at a glance; the rule does not rely on that.

Candidates were relaxed by the surrogate model for up to 100 steps (starting from iteration 10) and
evaluated with the DFT calculator below. Only structures from iteration $\ge 10$ were used in
the analysis, i.e.\ after relaxation had begun; earlier pre-relaxation placements were discarded.

\subsection{DFT settings}
\label{sec:dft}

Energies and forces were obtained with GPAW \cite{gpaw2014} in linear combination of atomic orbitals
(LCAO) mode (double-zeta polarised basis), the Perdew--Burke--Ernzerhof (PBE) exchange--correlation
functional \cite{pbe1996}, a $(1\times1\times1)$ k-point mesh,
Fermi--Dirac smearing (width 0.05~eV), and spin polarisation with Hund's rule coupling enabled; a
Pulay mixer and a fixed convergence criterion (energy $10^{-4}$~eV, density/eigenstates $10^{-3}$)
were used throughout. Note the deliberately modest computational setup (single-layer slab,
$\Gamma$-point sampling): the study targets trends across compositions, not converged
absolute energies (see Sec.~\ref{sec:scope}).

The interfacial Fe--O separation is a construction parameter rather than a computed
quantity: the reference films are built with the first oxygen layer of the substrate 2.3~\AA{} above
the metal plane, and relaxation preserves that spacing (2.300~\AA{} for the flat basin, i.e.\
unchanged to the printed precision, and 2.331~\AA{} for the island). Comparison with the 2.0--2.3~\AA{}
range reported for this interface by low-energy electron diffraction (LEED) analyses and earlier
calculations
\cite{urano1988,butler2001} therefore checks the model construction rather than testing the
electronic-structure method. Independent of this, the LCAO basis set is appreciably less complete
than a plane-wave or real-space representation \cite{larsen2009}, so absolute geometric parameters
are reported at fixed settings and are not interpreted as converged values.

\subsection{Flatness and energy metrics}
\label{sec:metrics}

Two quantities describe each sampled structure:

Film flatness $\dZ$---the vertical span of the metal film,
$\dZ = z(\text{metal})_\text{max} - z(\text{metal})_\text{min}$, taken over the film's
metal atoms (Fe; boron is a metalloid and is not part of the flatness metric).
$\dZ \approx 0$ is a flat, wetting film; large $\dZ$ is a 3D island (dewetted). A threshold of
$\dZ \le 1.0$~\AA{} separates the ``flat'' and ``island'' basins.

Relative energy $dE/N$---the per-atom energy above the system's global minimum,
$dE/N = (E - E_\text{min})/N$, where $E_\text{min}$ is the lowest energy found for that composition
and $N$ the atom count. By construction the global minimum is $0$~\eVA; per-atom normalisation makes
the two models directly comparable despite their differing atom counts (75 and 78 atoms).

For each model, the flat-basin ground state (lowest $dE/N$ among $\dZ \le 1.0$~\AA{}
structures) was compared against the overall minimum, isolating the energy penalty of the flat
configuration.

\subsection{Robustness of the method}
\label{sec:robustness}

Because the search relies on a biased-exploration scheme, three method parameters were varied on the
Fe/MgO model to test robustness (Supplementary Material): the rattle strength, the LCB parameter
$\kappa$, and the inclusion of a dipole correction. The principal conclusion was shown to be
insensitive to $\kappa$ and to the dipole correction, and to degrade only when the rattle strength
was reduced. This is the same model, method and reference layer as the Fe/MgO results of
Sec.~\ref{sec:results}, so the study bounds the method used for the two models reported here.

\subsection{Scope and limitations}
\label{sec:scope}

The structures retained here are not DFT-converged minima: candidates were relaxed with the
surrogate model and evaluated with a single DFT step, leaving residual forces of order
1--2~eV/\AA. Accordingly, ``lowest energy'' and ``basin'' denote the lowest DFT energy
\emph{encountered} in the search, not a converged minimum, and the results are interpreted
qualitatively (trends across compositions). The flat configuration is therefore described as a
\emph{higher-energy flat basin}, not as a metastable state---establishing metastability would require
converged relaxations, a Hessian (no imaginary modes) and a barrier between basins. A full
re-relaxation of representative structures is necessary for this purpose.

The models additionally build both phases on a body-centred cubic (bcc) Fe lattice, whereas the
experimentally reported structure of ultrathin Fe on MgO(001) is body-centred tetragonal
below about 10~\AA, changing to bcc only above that thickness \cite{urano1988}. The island branch
spans $\dZ \approx 1$--$6$~\AA, so the structures compared here lie in the thickness regime in which
the experimental film is not bcc. The flat--island comparison should therefore be read as a trend
obtained within a fixed lattice model, not as a prediction of absolute structural parameters.

One further property of the model bounds its interpretation: the strain convention is the
inverse of the experimental stack. The simulation cell takes the Fe lattice constant and the
substrate is compressed to match it, so the Fe film is unstrained in-plane, whereas in a
real junction a thin Fe film on bulk MgO absorbs the mismatch as in-plane strain and relieves it
through interfacial dislocations. This model therefore does not represent the strained-film
situation (see Sec.~\ref{sec:limitations}). The convention is the same for both models reported
here, so it does not enter the Fe/MgO-vs-Fe-B/MgO comparison at all.

Finally, the two models differ \emph{only} by the added boron: both are built on the same lattice
constant, the same substrate, the same reference-layer geometry and the same candidate-generation
schedule. What the comparison cannot separate is boron from \emph{any} added element---a second,
chemically different film composition would be needed for that, and none is analysed here.
